% Part: sets-functions-relations
% Chapter: sets
% Section: proofs-about-sets

\documentclass[../../../include/open-logic-section]{subfiles}

\begin{document}

\olfileid{sfr}{set}{prf}
\olsection{संचांबद्दलच्या सिद्धता}

\begin{editorial}
हा विभाग \verb|content/method/proofs| ने अधिक्रमित केला आहे आणि
या प्रकरणातून पूर्वनिर्धारितपणे काढला आहे.
\end{editorial}

\begin{explain}
आतापर्यंत ओळख करून दिलेले संच आणि संकेतलेखन
संच निर्दिष्ट करण्यासाठी व त्यांच्यातील संबंध व्यक्त
करण्यासाठी सोयीची लघुरूपे देतात. अशा संबंधांबद्दलचे
दावेही अनेकदा सिद्ध करावे लागतात. गणिती सिद्धतांची
ओळख नसेल, तर हे तुमच्यासाठी नवीन असेल. म्हणून एक
सोपे उदाहरण टप्प्याटप्प्याने पाहू. कोणत्याही संच~$X$
आणि $Y$ साठी $X \cap
(X \cup Y) = X$ हे नेहमी खरे असते, हे सिद्ध करू.
अशी संच-समानता कशी सिद्ध करायची? संच केवळ त्यांच्या
!!{element}s वरून ठरतात, म्हणजे त्यांत तेच !!{element}s
असतील तेव्हाच ते समान असतात, हे आठवा. म्हणून येथे
(a) $X \cap (X \cup
Y)$ मधील प्रत्येक !!{element} हा $X$ चाही !!{element}
आहे आणि, उलट, (b) $X$ मधील प्रत्येक !!{element}
हा $X \cap (X \cup Y)$ चाही !!{element} आहे, हे
सिद्ध करायचे आहे. दुसऱ्या शब्दांत, (a) $X \cap (X \cup Y) \subseteq X$
आणि (b) $X \subseteq X \cap (X \cup Y)$ हे दोन्ही दाखवतो.

“$X \cap (X
\cup Y)$ मधील प्रत्येक !!{element}~$z$ हा $X$ चाही !!{element}
आहे” असा सर्वसाधारण दावा सिद्ध करताना प्रथम कोणताही
$z \in X \cap (X \cup Y)$ दिला आहे असे मानतो आणि या
गृहीतकापासून $z \in X$ सिद्ध करतो. ही पद्धत “सर्वसाधारण
सोपाधिक सिद्धता” म्हणून परिचित असेल. या सिद्धतेत
गृहीतक आणि निष्कर्ष यांतील व्याख्यांचाही उपयोग करावा
लागेल; येथे उदाहरणार्थ “$\cap$” आणि “$\cup$” यांच्या
व्याख्यांचा. म्हणून (a) साठीचा युक्तिवाद असा:


\begin{quote}
(a) प्रथम $X \cap (X \cup Y) \subseteq X$ दाखवायचे आहे;
म्हणजे $\subseteq$ च्या व्याख्येनुसार कोणत्याही~$z$ साठी
$z \in X \cap (X \cup Y)$ असेल, तर $z \in X$ हे दाखवायचे.
म्हणून $z \in X \cap (X \cup
Y)$ असे माना. $z$ हा दोन संचांच्या छेदाचा !!{element}
तेव्हाच आणि तेव्हाच असतो जेव्हा तो दोन्ही संचांचा
!!{element} असतो. म्हणून $z \in X$ आणि $z \in X \cup Y$
असे निष्पन्न होते. विशेषतः $z \in X$, हेच दाखवायचे होते.
\end{quote}

याने सिद्धतेचा पहिला भाग पूर्ण झाला. शेवटच्या
पायरीत आपण एक तथ्य वापरले: एखाद्या गृहीतकापासून
संयुती ($z \in X$ आणि $z \in X
\cup Y$) निष्पन्न झाली, तर तिचा प्रत्येक संयुतीघटक
त्याच गृहीतकापासून निष्पन्न होतो. हा नियम “संयुती-निरसन”
किंवा $\land$Elim म्हणून परिचित असेल. आता (b) सिद्ध करू:

\begin{quote}
(b) आता $X \subseteq X \cap (X \cup Y)$ सिद्ध करू; म्हणजे
$\subseteq$ च्या व्याख्येनुसार कोणत्याही~$z$ साठी
$z \in X$ असेल, तर $z \in X \cap
(X \cup Y)$ हे दाखवू. $z \in X$ असे माना. $z \in X
\cap (X \cup Y)$ दाखवण्यासाठी “$\cap$” च्या व्याख्येनुसार
(i) $z \in X$ आणि (ii) $z \in X \cup Y$ हे दाखवावे
लागेल. (i) हे आपले गृहीतकच आहे; म्हणून अधिक
सिद्ध करण्याची गरज नाही. (ii) साठी आठवा की $z$
हा संचांच्या संयोगाचा !!{element} तेव्हाच आणि तेव्हाच
असतो जेव्हा तो त्या संचांपैकी किमान एकाचा !!{element}
असतो. $z \in X$ आणि $X \cup Y$ हा $X$ व $Y$ यांचा
संयोग असल्याने येथे ही अट पूर्ण होते. म्हणून
$z \in X \cup Y$. आपण (i) $z \in X$ आणि
(ii) $z \in X \cup Y$ ही दोन्ही दाखवली; त्यामुळे
“$\cap$” च्या व्याख्येनुसार $z \in X \cap (X \cup Y)$.
\end{quote}

हा युक्तिवाद काहीसा लांबलचक होता; पण संच आणि
त्यांच्यातील संबंधांविषयी आपण कसा विचार करतो हे
तो दाखवतो. सहसा आपण इतक्या स्पष्टपणे सर्व काही
लिहीत नाही; विशेषतः प्रत्येक व्याख्या पुन्हा सांगत
नाही. अधिक प्रगत पाठ्यात याच निकालाची सिद्धता
खूप संक्षिप्त असेल. ती पुढीलप्रमाणे दिसू शकते.
\end{explain}

\begin{prop}[शोषण]
सर्व संच $X$, $Y$ साठी,
\[
X \cap (X \cup Y) = X
\]
\end{prop}

\begin{proof}
(a) $z \in X \cap (X \cup Y)$ असे माना. मग $z \in X$;
म्हणून $X \cap (X
  \cup Y) \subseteq X$.

(b) आता $z \in X$ असे माना. मग $z \in X \cup Y$ सुद्धा;
त्यामुळे $z \in X \cap (X \cup Y)$ सुद्धा. म्हणून
$X \subseteq X \cap (X \cup
  Y)$.
\end{proof}

\begin{prob}
$X \cup (X \cap Y) = X$ हे तपशीलवार सिद्ध करा.
मग त्याची संक्षिप्त सिद्धता द्या. (सूचना:
$X \cup (X \cap Y) \subseteq X$ या दिशेसाठी
प्रकरणांनुसार सिद्धता, म्हणजे $\lor$Elim, लागेल.)
\end{prob}

\end{document}
